% =====================================================================
%  BOZZA - Capitolo 2 (italiano)
%  Contributo: consapevolezza multihop e protocolli di beaconing
%
%  Richiede i pacchetti gia' presenti nel preambolo del wrapper:
%  amsmath, amssymb, siunitx, enumitem, booktabs.
%  Le macro seguenti sono ridefinite con \providecommand per sicurezza.
% =====================================================================
\providecommand{\bimin}{BI_{\min}}
\providecommand{\bimax}{BI_{\max}}
\providecommand{\nthr}{N_{\mathrm{thr}}}

\chapter{Estensione della consapevolezza: meccanismi multihop e protocolli di beaconing}
\label{ch:contributo}

Questo capitolo descrive il contributo principale della tesi apportato al simulatore: le due
modalita' di multihop \emph{append} e \emph{forwarded}, che estendono la consapevolezza dei nodi
oltre il singolo hop.

Prima pero' si richiama il modello di comunicazione su cui poggia il simulatore, ovvero la
struttura del beacon e i protocolli che ne regolano gli intervalli di trasmissione (SBP, ADAB,
ACAB).

Il capitolo si chiude con alcuni possibili miglioramenti progettati per contenere il costo delle
due modalita' multihop, ma non ancora attivi nell'implementazione corrente.

% ---------------------------------------------------------------------
\section{Modello di comunicazione e struttura del beacon}
\label{sec:commodel}
ProSafe esegue la scoperta di prossimit\`a a livello MAC IEEE~802.11, senza connettersi ad access
point o utilizzare servizi di livello IP. I beacon sono brevi pacchetti inviati in broadcast.
Il payload modellato nel simulatore (nella classe
\texttt{Beacon}) contiene:
\begin{itemize}[leftmargin=1.4em,nosep]
  \item \textbf{sender id} --- identificatore univoco della boa trasmittente ($16$ byte);
  % \item \textbf{flag di mobilit\`a} --- indica se il nodo \`e mobile ($1$ byte);
  \item \textbf{posizione} --- coppia di coordinate correnti relative al piano di simulazione ($8$ byte);
  \item \textbf{timestamp} --- istante temporale di trasmissione ($4$ byte);
  \item \textbf{lista dei nodi conosciuti} --- ciascuna voce e' costituita da id, timestamp,
  posizione e hop del nodo vicino, le prime tre componenti hanno un peso di $28$ byte, mentre la distanza in hop \`e usata solo internamente per le operazioni del simulatore ed e' trascurata;
  \item \textbf{origin id} e \textbf{hop limit} --- $16+4$ byte, valorizzati solamente in
  modalit\`a forwarded.
\end{itemize}
Detto $k$ il numero di nodi annunciati nella lista, la dimensione del frame in byte
\`e
\begin{equation}
  S(k) = 33 + 28\,k \,\,(+ 20\,_{\text{forwarded}})
  \label{eq:beacon-size}
\end{equation}
Il frame base \`e quindi di $33$ byte e cresce di $28$ byte per ogni nodo
annunciato. La dimensione del beacon quindi cresce in proporzione al numero di nodi
che la boa conosce.

L'accesso al canale e' mediato dai principi della DCF (Distributed coordination function): il nodo rileva il mezzo, attende un DIFS, calcola il tempo di backoff e trasmette, senza ricorrere a RTS/CTS o ACK.
La dinamicita' della rete varia a seconda dello \texttt{scenario} impostato (static, ramp, random), questo parametro va ad impattare sulla variazione del numero di nodi durante la durata della simulazione.
Per i fini della valutazione dei risultati ci si e' concentrati sullo scenario random, nel quale il numero dei nodi sale e decresce periodicamente di un fattore percentuale stabilito dalla configurazione; allo scopo di simulare il comportamento dinamico di una rete reale, dove i nodi entrano ed escono dal raggio radio in continuazione.
Ogni nodo mantiene in memoria una tabella dei vicini, aggiornata a ogni ricezione e svuotata dalle informazioni obsolete periodicamente.

Oltre alla struttura del beacon, il modello di comunicazione comprende i protocolli che decidono
quando una boa trasmette. Sono tre: SBP, ADAB e ACAB; agiscono solo sull'istante di
generazione del beacon e lasciano invariato il CSMA/CA sottostante. Questi protocolli erano preesistenti nel simulatore e 
vengono richiamati qui solo come contesto perche' ogni modalit\`a multihop introdotta pi\`u
avanti pu\`o essere abbinata a uno qualunque di essi.

\subsection{Static Beaconing Protocol (SBP)}
\label{sec:sbp}
Lo SBP e' il protocollo di riferimento. Ogni boa trasmette un beacon a intervallo
fisso $BI_{\text{static}}$, indipendentemente dalle condizioni di rete, ascolta in
continuazione e aggiorna la tabella dei vicini a ogni ricezione. \`E semplice e
prevedibile, ma in scenari densi le collisioni riducono l'affidabilita' mentre in scenari sparsi l'intervallo fisso porta a
trasmissioni inutilmente frequenti. Questi limiti motivano il bisogno di protocolli adattivi.

\subsection{Adaptive Density-Aware Beaconing (ADAB)}
\label{sec:adab}
ADAB regola l'intervallo in base a una sola grandezza locale: il numero di vicini
in tabella. A ogni decisione il nodo conta i vicini correnti $N_t$ e con essi calcola un fattore di densita' normalizzato
\begin{equation}
  f_t = \min\!\left(1,\; \frac{N_t}{\nthr}\right),
  \qquad \nthr = 15,
  \label{eq:adab-density}
\end{equation}
dove $\nthr$ \`e la soglia oltre la quale la congestione diventa rilevante.
L'intervallo successivo \`e una funzione quadratica del fattore di densit\`a,
interpolata fra un minimo e un massimo, con un piccolo jitter moltiplicativo che
desincronizza i nodi:
\begin{equation}
  BI_{t+1} = \Big[\,\bimin + f_t^{\,2}\,(\bimax-\bimin)\,\Big]\,(1+\epsilon_t),
  \qquad \epsilon_t \sim \mathcal{U}[-\eta,\eta],\;\; \eta=0.05,
  \label{eq:adab}
\end{equation}
con $BI_{t+1}$ vincolato all'intervallo $[\bimin,\bimax]$. Con pochi vicini
$f_t\approx 0$ e la boa trasmette vicino a $\bimin$, per una scoperta rapida; man
mano che $N_t$ si avvicina a $\nthr$ il termine quadratico porta l'intervallo verso
$\bimax$, riducendo le trasmissioni dove la contesa \`e maggiore. La forma
quadratica tiene il nodo reattivo a bassa densit\`a e ne accelera il rallentamento
ad alta densit\`a. I valori usati sono $\bimax = \SI{5.0}{\second}$ e $\bimin$
configurabile (negli esperimenti $\bimin \in \{0.25, 0.5, 1.0\}\,\si{\second}$).

\subsection{Adaptive Context-Aware Beaconing (ACAB)}
\label{sec:acab}
La densit\`a non \`e l'unico segnale locale utile a determinare se un beacon vada
trasmesso oppure no. ACAB combina tre segnali: la densit\`a dei vicini, il tempo trascorso
dall'ultimo contatto e la mobilit\`a del nodo.

\paragraph{Componente di densit\`a.} Come in ADAB, si calcola un conteggio
normalizzato dei vicini, con una soglia pi\`u bassa:
\begin{equation}
  s_{\mathrm{den}} = \min\!\left(1,\;\frac{N_t}{\nthr^{\mathrm{ACAB}}}\right),
  \qquad \nthr^{\mathrm{ACAB}} = 10 .
  \label{eq:acab-den}
\end{equation}

\paragraph{Componente di tempo dall'ultimo contatto.} Se la boa ha ricevuto un beacon da poco
non ha molta fretta di pubblicizzarsi; se invece e' rimasta a lungo in silenzio, conviene che
trasmetta. Detto $\Delta = t - t_{\text{ultimo contatto}}$ il tempo passato dall'ultimo beacon
ricevuto e $\tau_c = \SI{20}{\second}$ la durata oltre la quale un contatto non e' piu' considerato
recente:
\begin{equation}
  s_{\mathrm{con}} =
  \begin{cases}
    \max\!\big(0,\; 1 - \Delta/\tau_c\big) & N_t>0 \text{ e contatto registrato},\\
    0 & \text{altrimenti.}
  \end{cases}
  \label{eq:acab-con}
\end{equation}
Un nodo contattato da poco d\`a $s_{\mathrm{con}}\to 1$ (rallenta); un nodo isolato
o silenzioso da tempo d\`a $0$ (trasmette pi\`u spesso).

\paragraph{Componente di mobilit\`a.} Un nodo veloce cambia posizione rapidamente e
dovrebbe aggiornarla pi\`u spesso. Con velocit\`a $v=\lVert\mathbf{v}\rVert$ e una
velocit\`a di riferimento $v_{\text{ref}}$ (la velocit\`a di default, qui
$\SI{15}{\metre\per\second}$):
\begin{equation}
  s_{\mathrm{mob}} = \min\!\left(1,\;\frac{v}{v_{\text{ref}}}\right).
  \label{eq:acab-mob}
\end{equation}

\paragraph{Combinazione e calcolo dell'intervallo.} I tre punteggi sono combinati
con pesi $(w_d,w_c,w_m)=(0.4,0.3,0.3)$ che sommano a uno. La mobilit\`a riduce il
backoff, da cui il termine $1-s_{\mathrm{mob}}$:
\begin{equation}
  s = w_d\,s_{\mathrm{den}} + w_c\,s_{\mathrm{con}} + w_m\,(1 - s_{\mathrm{mob}}).
  \label{eq:acab-blend}
\end{equation}
Il punteggio $s$ viene poi usato al posto di $f_t$ nel mapping quadratico con
jitter di ADAB (Eq.~\ref{eq:adab}). Cos\`i ACAB rallenta una boa densa, contattata
di recente e ferma, e tiene pi\`u frequente una boa isolata, silenziosa o veloce.
Rispetto ad ADAB, che usa la sola densit\`a, ACAB costruisce il punteggio di
backoff da tre segnali; il costo \`e un po' di stato in pi\`u (velocit\`a e istante
dell'ultimo contatto), trascurabile sui dispositivi considerati. Il calcolo
dell'intervallo \`e dettagliato nel Capitolo~\ref{ch:simulatore}.

% ---------------------------------------------------------------------
\section{Estensione della consapevolezza dei nodi}
\label{sec:multihop}
Ad un nodo interessa conoscere la posizione e lo stato dei nodi che gli stanno attorno nella
rete. Gia' a un solo hop un beacon non si limita ad annunciare il proprio mittente: porta con se'
la lista dei nodi che il mittente conosce, e cosi' chi lo riceve viene a sapere anche di nodi che
non sente in modo diretto.

Il contributo di questa tesi riguarda invece l'estensione della consapevolezza oltre il singolo
hop. Spesso a un nodo serve sapere anche dei nodi piu' lontani, o coperti da altri nodi, e
propagare piu' in fretta le informazioni in suo possesso; ma ogni trasmissione spesa per
propagarle compete per lo stesso canale, gia' sotto pressione e privo di conferme. I due
meccanismi seguenti sono due risposte opposte a questo compromesso.

\subsection{Append: estensione passiva}
\label{sec:append}
In modalit\`a append una boa allega, a ogni beacon, i nodi di cui \`e a
conoscenza: sia i vicini diretti (a $1$ hop) sia i nodi scoperti
in modo indiretto dalle liste annunciate dagli altri. Non si manda nessun frame in
pi\`u: la consapevolezza viaggia sui beacon che il nodo invierebbe comunque, e l'unico costo
\`e nella dimensione del frame, che cresce di $28$ byte per ogni nodo annunciato
(Eq.~\ref{eq:beacon-size}).

Alla ricezione di un beacon, la boa aggiorna prima i vicini diretti con il
mittente, poi integra la lista annunciata in una tabella di \emph{nodi scoperti}.
Per ogni voce $(j, t_j, p_j, h_j)$ della lista, con $h_j$ la distanza in hop dal
mittente del beacon, il nodo ricevente registra la stessa voce a distanza
\begin{equation}
  h_j' = h_j + 1,
  \label{eq:append-hop}
\end{equation}
perch\'e quel nodo si trova un hop pi\`u lontano da chi riceve che da chi ha
inviato. La voce viene scartata se $j$ \`e il nodo stesso o se \`e gi\`a un vicino
diretto (i nodi scoperti non duplicano i vicini diretti). Negli altri casi la
tabella dei nodi scoperti viene aggiornata solo se il beacon porta informazione
pi\`u fresca, cio\`e se
\begin{equation}
  t_j > t_j^{\text{mem}},
  \label{eq:append-fresh}
\end{equation}
dove $t_j^{\text{mem}}$ \`e il timestamp dell'ultima informazione memorizzata su
$j$ (oppure $-\infty$ se $j$ \`e nuovo). Cos\`i ogni nodo tende a conservare la
posizione pi\`u recente nota di ciascun nodo raggiungibile.

\subsection{Forwarded: estensione attiva}
\label{sec:forwarded}
In modalit\`a forwarded una boa ritrasmette i beacon che riceve. Ogni beacon porta
un \emph{origin id} (il mittente originale) e un \emph{hop limit} $\ell$, un
contatore in stile TTL. Quando una boa riceve un beacon che non ha originato e con
$\ell > 0$, pu\`o accodare una copia da rilanciare. Al rilancio il contatore viene
decrementato,
\begin{equation}
  \ell' = \ell - 1,
  \label{eq:ttl}
\end{equation}
mentre origine, timestamp e payload restano invariati; la boa che rilancia diventa
il mittente ai fini del calcolo del range, ma non altera il contenuto. Mantenere
origine e timestamp permette di riconoscere i duplicati e di evitare che un nodo
rilanci un beacon da s\'e originato. I rilanci passano per la stessa pipeline CSMA
dei beacon propri, dopo un piccolo jitter casuale
$\delta \sim \mathcal{U}[0, \SI{0.05}{\second}]$ che desincronizza i molti
ricevitori dello stesso frame, i quali altrimenti contenderebbero il canale tutti
insieme.

Senza controlli il forwarding attivo pu\`o generare un \emph{broadcast storm}.
L'implementazione adotta tre meccanismi di contenimento.
\begin{itemize}[leftmargin=1.4em,nosep]
  \item \textbf{Deduplicazione per origine.} Ogni boa memorizza, per ciascun
  origin id $o$, il timestamp pi\`u recente gi\`a inoltrato $\tau_o$. Un beacon di
  origine $o$ e timestamp $t$ viene accodato solo se
  \begin{equation}
    t > \tau_o,
    \label{eq:dedup}
  \end{equation}
  cio\`e solo se pi\`u fresco dell'ultimo gi\`a inoltrato da quella origine. Questo
  evita di rilanciare pi\`u volte la stessa informazione e spezza i cicli.
  \item \textbf{Limite del TTL.} Il contatore $\ell$, decrementato a ogni rilancio
  (Eq.~\ref{eq:ttl}), limita la profondit\`a di propagazione; negli esperimenti il
  valore tipico \`e $1$ hop.
  \item \textbf{Coda limitata.} I rilanci pendenti stanno in una coda di dimensione
  massima \texttt{pending\_queue\_limit}. Quando \`e piena i nuovi beacon vengono
  scartati senza essere registrati, cos\`i che una copia successiva possa entrare
  se nel frattempo si libera spazio.
\end{itemize}
In sintesi, la modalit\`a forwarded corrente accoda e rilancia ogni beacon fresco,
entro TTL e non proprio, con il solo limite della coda.

\subsection{Copertura e costo di canale}
\label{sec:tradeoff}
I due meccanismi sono scelte opposte. L'append \`e economico, perch\'e non aggiunge
nessun frame, ma \`e limitato dalla freschezza dell'informazione e dalla dimensione
crescente del beacon (Eq.~\ref{eq:beacon-size}): pi\`u nodi conosce una boa,
pi\`u grande, e pi\`u costoso in airtime, diventa ogni suo beacon. Il forwarded
estende la portata della consapevolezza, perch\'e un beacon pu\`o
raggiungere nodi fuori dal range del mittente originale, ma aumenta il traffico, le
collisioni e la latenza. Quale dei due convenga in una rete broadcast densa, a
dominio unico e senza conferme \`e valutato nel capitolo dedicato ai risultati.

% ---------------------------------------------------------------------
\section{Possibili miglioramenti}
\label{sec:improvements}
Oltre ai due meccanismi multihop appena descritti, durante lo sviluppo sono stati progettati due
accorgimenti per contenerne il costo: un limite di hop per l'append e un gate di forwarding
density-aware. Entrambi sono presenti nel codice ma non attivi negli esperimenti di questa tesi,
e la loro valutazione \`e lasciata ai lavori futuri.

\subsection{Limite di hop per l'append}
\label{sec:append-hoplimit}
In modalit\`a append (Sezione~\ref{sec:append}) una boa annuncia tutti i nodi scoperti, a
prescindere da quanto sono lontani in hop; in reti dense questo allunga la lista annunciata e
quindi la dimensione del beacon (Eq.~\ref{eq:beacon-size}). Un eventuale limite di hop per
l'append (\texttt{append\_hop\_limit}) \`e gi\`a caricato dalla configurazione ma \emph{non}
applicato. La distanza $h_j'$ di ogni voce (Eq.~\ref{eq:append-hop}) resta comunque registrata, e
basterebbe usarla per non annunciare i nodi oltre una certa distanza in hop, tenendo cos\`i sotto
controllo la crescita del beacon. La taratura del limite e l'effetto su copertura e airtime sono
lasciati ai lavori futuri.

\subsection{Gate di forwarding density-aware}
\label{sec:fwdgate}
La modalit\`a forwarded della Sezione~\ref{sec:forwarded} rilancia ogni beacon
fresco, entro TTL e non proprio, con il solo limite della coda. In reti dense
questo pu\`o produrre molti rilanci ridondanti, perch\'e molte boe nello stesso
intorno rilanciano lo stesso beacon. Un'estensione progettata, ma non attiva
nell'implementazione usata in questa tesi, \`e un \emph{gate di forwarding} che
decide, una sola volta al momento dell'accodamento, se rilanciare un beacon
ricevuto, con probabilit\`a
\begin{equation}
  P_{\text{fwd}} =
  \begin{cases}
    1 & N \le n_0,\\[2pt]
    \dfrac{n_0}{N}\cdot \max\!\big(0.3,\;1-f_q\big) & N > n_0,
  \end{cases}
  \label{eq:fwdgate}
\end{equation}
dove $N$ \`e il numero di vicini del nodo che rilancia, $n_0$ \`e un numero base
atteso di inoltranti per cascata ($\texttt{forward\_density\_baseline}=5$) e $f_q$
\`e il punteggio di backoff al quadrato calcolato per ultimo dallo scheduler. Il
fattore $n_0/N$ dirada i rilanci nelle zone dense, cos\`i che in media solo
$\approx n_0$ degli $N$ candidati rilancino. Il fattore $\max(0.3,\,1-f_q)$ riusa
lo stesso segnale di congestione locale che governa la frequenza dei beacon propri:
una boa che sta gi\`a rallentando i propri beacon riduce anche i rilanci, fino a un
minimo di $0.3$.

Nel codice il gate \`e presente ma commentato (\texttt{should\_forward}), e il
parametro \texttt{forward\_density\_baseline} \`e caricato ma inattivo: gli
esperimenti usano la modalit\`a forwarded senza gate. La sua valutazione, cio\`e
l'impatto su rilanci, collisioni, latenza e copertura e la taratura di $n_0$, \`e
lasciata ai lavori futuri.

% ---------------------------------------------------------------------
\section*{Sintesi del capitolo}
Questo capitolo ha presentato il contributo della tesi. I due meccanismi multihop
estendono la conoscenza dei nodi in modo opposto: l'append (passivo) aggiunge ai
beacon la lista dei nodi conosciuti senza traffico extra, incrementando la distanza
in hop e tenendo l'informazione pi\`u fresca; il forwarded (attivo) rilancia i
beacon entro un TTL, con deduplicazione per origine e coda limitata a contenere il
broadcast storm. I tre protocolli di beaconing (SBP, ADAB, ACAB), gi\`a presenti nel
simulatore e parte del modello di comunicazione, decidono quando
trasmettere e si possono combinare con entrambe le modalit\`a. Infine si sono
descritti due possibili miglioramenti, un limite di hop per l'append e un gate di
forwarding density-aware, progettati ma non ancora attivi. Il prossimo capitolo
descrive come tutto questo \`e realizzato nel simulatore.
